\documentclass[11pt,a4paper]{ltjsarticle}

% ── Packages ──
\usepackage[margin=2.5cm]{geometry}
\usepackage{luatexja-fontspec}
\setmainjfont{IPAexMincho}
\setsansjfont{IPAexGothic}
\usepackage{amsmath,amssymb}
\usepackage{booktabs}
\usepackage{longtable}
\usepackage{array}
\usepackage{graphicx}
\usepackage[dvipsnames]{xcolor}
\usepackage[numbers,sort&compress]{natbib}
\usepackage{hyperref}
\hypersetup{colorlinks=true,linkcolor=blue,citecolor=blue,urlcolor=blue}
\usepackage{enumitem}
\setlist{nosep}

% ── Layout tweaks ──
\setlength{\parindent}{0pt}
\setlength{\parskip}{0.6em}
\renewcommand{\baselinestretch}{1.15}

% ── Custom commands ──
\newcommand{\HRule}{\rule{\linewidth}{0.4pt}}

% ══════════════════════════════════════════════════════════════════
\begin{document}

% ── Title ──
\begin{center}
{\LARGE\bfseries タンパク質配列からの階層的糖鎖予測により\\
構造クラスレベルでの決定境界が明らかに}

\vspace{0.5em}
{\large Hierarchical glycan prediction from protein sequence\\
reveals a determination boundary at the structural-class level}

\vspace{1em}
{\large 松井 祐介$^{1,2,*}$}\\[0.3em]
{\normalsize $^1$名古屋大学大学院医学系研究科 総合保健学専攻 生体情報・医工計測学分野\\
$^2$名古屋大学 糖鎖生命コア研究所（iGCORE）}\\[0.3em]
{\normalsize $^*$責任著者：松井 祐介}
\end{center}
\HRule
\vspace{0.5em}

% ══════════════════════════════════════════════════════════════════
\section*{要旨}

\textbf{動機：}
糖鎖修飾は最も複雑な翻訳後修飾の一つであるが、特定のタンパク質部位にどの糖鎖構造が付加されるかを予測することは未解決の課題である。糖鎖生物学データは互換性のないスキーマを持つ専門データベースに分散しており、糖鎖のアイデンティティに対するタンパク質配列と細胞環境のそれぞれの寄与は定量化されていない。

\textbf{結果：}
我々は、6つの公開データベースから78,263ノード（10エンティティ型）と250万エッジ（14関係型）にわたる統合ナレッジグラフglycoMusubiを構築した。glycoMusubiを用いて、糖鎖予測を特異性が増加する6つのタスクに分解する階層的ベンチマークを確立した：N結合型部位分類（F1 = 0.924）、部位ランキング（Recall@3 = 0.683）、糖鎖修飾タイプ予測（精度 = 0.844）、構造ファミリー分類（Top-2精度 = 0.858）、構造クラスタ予測（Hits@5 = 0.786、20クラスタ）、正確な糖鎖同定（MRR = 0.066--0.118、2,357--3,345候補）。クラスタレベルとIDレベルの性能間の約10倍のギャップにより、定量的な決定境界が明らかになった：タンパク質配列は構造クラスを確実に予測するが、正確な糖鎖のアイデンティティには追加情報が必要である。シングルセル酵素発現データ（Tabula Sapiens; 748細胞型、339酵素）の統合は、頻度ベースラインを超える予測改善をもたらさなかった（MRR = 0.035 対 0.105）。これは現在の酵素-糖鎖アノテーションがN結合型糖鎖のわずか39\%しかカバーしていないためである。これらの結果は糖鎖修飾決定の2つの層を描出し、アノテーションの疎さが主要なボトルネックであることを特定した。

\textbf{利用可能性と実装：}
glycoMusubiソフトウェアとベンチマークスイートは
\url{https://github.com/ymatts/glycoMusubi}
にてMITライセンスで公開している。ナレッジグラフと学習済みモデルはZenodoに登録予定である。

\textbf{連絡先：}松井 祐介

\textbf{補足情報：}補足データは\textit{Bioinformatics}オンラインで入手可能。

\HRule

% ══════════════════════════════════════════════════════════════════
\section{序論}

糖鎖---タンパク質や脂質に共有結合する複雑な糖質---は、免疫認識、細胞シグナル伝達から病原体侵入、がん進行に至るまで、事実上すべての細胞表面の生物学的プロセスを媒介する~\cite{varki2017,reily2019}。数十年にわたる構造解析にもかかわらず、計算科学的な根本的問題は未解決のままである：タンパク質配列と細胞コンテキストが与えられたとき、特定の部位にどの糖鎖構造が付加されるのか？

この問題が困難なのは、糖鎖修飾が複数のレベルで決定されるためである。タンパク質配列と局所的な部位コンテキストは、どの糖転移酵素が基質にアクセスできるかに影響し、可能な糖鎖の構造クラスを制約する~\cite{schjoldager2020,narimatsu2019}。しかし、正確な糖鎖構造はさらに、細胞内で発現する糖鎖修飾酵素の組み合わせ、ヌクレオチド糖ドナーの利用可能性、および競合する生合成経路の反応速度論に依存する~\cite{stanley2022,varki2022essentials}。これらの「配列決定」因子と「環境決定」因子のそれぞれの寄与は定性的に議論されてきた~\cite{schjoldager2020}が、体系的に定量化されたことはない。

近年の計算的アプローチはこの問題の一部に取り組んできた。InSaNNE~\cite{kellman2024}はリカレントニューラルネットワークを訓練し、タンパク質配列特徴量から糖鎖構造を予測し、57糖鎖付加部位にわたる199糖鎖のキュレーションされたセットでAUC = 0.92を達成した。しかし、このアプローチは検索ランキングではなく二値の有無判定を評価し、制限された候補セットを使用し、予測難易度の階層を分解していない。BojarとLisacek~\cite{bojar2022}は糖鎖インフォマティクスと機械学習のアプローチをレビューし、糖鎖特性予測とタンパク質コンテキストからの糖鎖予測という逆問題とのギャップを指摘した。CoNglyPred~\cite{congly2025}はESM-2とAlphaFold2の特徴量をN結合型部位予測に組み合わせ、我々のTask~1を補完する。GlycanML~\cite{glyml2024}は11の糖鎖特性予測タスク（免疫原性、分類学、組織発現）のベンチマークを確立したが、これらは糖鎖構造\textit{から}予測するものであり、タンパク質コンテキスト\textit{から}糖鎖構造を予測するものではない。

生物学的側面では、Yadavら~\cite{yadav2022}がゴルジ体での糖鎖プロセシングを最適情報符号化チャネルとしてモデル化し、構造多様性が酵素反応速度論によって制約されることを示した---これは我々の経験的な決定境界と整合する理論的枠組みである。Chrysinasら~\cite{chrysinas2024}は我々と同じTabula Sapiensアトラスから糖鎖修飾酵素の発現を解析し、末端プロセシング酵素が高度に細胞型特異的である一方、コア経路酵素はユビキタスに発現していることを見出した。彼らの定性的観察は、酵素発現だけでは正確な糖鎖のアイデンティティを解決できないという我々の定量的知見と一致する。Glycoimpact~\cite{glycoimpact2024}は糖鎖制約のためのアミノ酸置換スコアリングマトリックスを定義し、配列が糖鎖を制約するが決定はしないという「制約付き生合成」の概念を導入した。

データの断片化も進歩を阻んできた。糖鎖構造はGlyTouCan~\cite{tiemeyer2017}に、タンパク質アノテーションはUniProt~\cite{uniprot2023}に、酵素-糖鎖の生合成関係はGlyGen~\cite{york2020}に、阻害剤活性はChEMBL~\cite{zdrazil2024}に、部位特異的糖鎖割り当てはGlyConnect~\cite{alocci2019}に存在する。各データベースは異なるスキーマと識別子を使用しており、統合解析を困難にしている。糖鎖表現における最近の進歩---グラフニューラルネットワークエンコーディングのためのSweetNet~\cite{burkholz2022}、高次表現のためのGIFFLAR~\cite{thomes2024}、およびモチーフベースの解析のためのglycowork~\cite{thomes2021}---は糖鎖特性予測を改善したが、タンパク質コンテキストからの糖鎖予測という逆問題には取り組んでおらず、決定境界も定量化していない。

本研究では3つの貢献を行う。第一に、6つの公開データベースから自動化された再現可能なETLパイプラインを通じて78,263ノード（10型）と2,508,960エッジ（14関係型）を統合するヘテロジニアスナレッジグラフ\textbf{glycoMusubi}を構築した。第二に、二値のN結合型部位分類から2,357候補にわたる正確な糖鎖同定まで、特異性が増加する6つのタスクからなる\textbf{階層的予測ベンチマーク}を確立した。第三に、シングルセル酵素発現データを用いた\textbf{細胞コンテキストの体系的調査}を実施し、現在の酵素-糖鎖アノテーションが決定ギャップを橋渡しするには不十分であるという最初の定量的証拠を提供した。これらの解析を合わせて、配列予測可能な構造クラスと環境依存の糖鎖アイデンティティの間の境界を描出し、分野を前進させるために埋めるべき具体的なデータギャップを特定した。

% ══════════════════════════════════════════════════════════════════
\section{材料と方法}

\subsection{ナレッジグラフの構築}

glycoMusubi ETLパイプラインは4つの自動化ステージで構成される：\textbf{ダウンロード}、\textbf{クリーニング}、\textbf{ビルド}、\textbf{バリデーション}（補足図~S1）。データは、設定可能なリトライロジックと指数バックオフを備えたレート制限付きAPIクライアントを通じて6つの公開データベースから取得される。

\textbf{データソース。}
GlyGenは糖転移酵素アノテーションと糖鎖マスターリストを提供する。GlyTouCanはGlyCosmos SPARQList APIを通じてWURCS表記の糖鎖構造を提供する。UniProtは疾患関連、遺伝的バリアント、糖鎖付加部位の特徴を含むタンパク質アノテーションを提供する。ChEMBLは糖転移酵素阻害剤の活性データを提供する。PTMCodeは翻訳後修飾の部位レベルクロストークネットワークデータを提供する。GlyConnectはタンパク質と位置アノテーション付きの部位特異的糖鎖割り当てを提供する。

\textbf{スキーマの適用。}
KGスキーマは10のノード型と14の関係型を定義する（補足表~S1）。ノード識別子はソース固有のパターンに対して検証される（GlyTouCan: \texttt{\^{}G\textbackslash d\{5\}[A-Z]\{2\}\$}; UniProt: Swiss-Protアクセッション形式; ChEMBL: \texttt{\^{}CHEMBL\textbackslash d+\$}）。エッジは宣言されたソース-ターゲット型ペアに制約される。バリデーションチェックには、孤立ノード検出、重複エッジ除去、自己ループ排除、非巡回関係における循環検出が含まれる。

\textbf{糖鎖付加部位の処理。}
UniProt糖鎖付加部位は、実験的に検証されたエントリのみ保持するようフィルタリングされる（エビデンスコードECO:0000269, ECO:0000303, ECO:0000250）。予測および低信頼度のアノテーションは除外される。部位位置は標準化されたフォーマット\texttt{SITE::<UniProtID>::<Position>::<Residue>}を使用してカノニカルアイソフォーム座標に正規化される。

\textbf{PTMクロストークの統合。}
PTMCodeクロストークエッジ（最小信頼度スコア0.30）は、6つのPTM型（リン酸化、メチル化、アセチル化、ユビキチン化、SUMO化、糖鎖修飾）にわたる機能的相互作用のエビデンスを持つ修飾部位のペアをリンクする。

\subsection{ナレッジグラフ埋め込みモデル}

構造クラスタリングと検索のための糖鎖表現を得るため、glycoMusubiグラフ上でナレッジグラフ埋め込みモデル（GlycoKGNet）を訓練した。

\textbf{モダリティ特異的エンコーダ。}
糖鎖ノードはGlycanTreeEncoderによりエンコードされる。これはWURCSパースされた分岐木構造に対してボトムアップのツリーメッセージパッシング層（子アテンションとGRU状態更新付き）を適用し、続いてトップダウンの精緻化とアテンションプーリングを行い、256次元の糖鎖埋め込みを生成する。タンパク質ノードは、事前計算されたESM-2~\cite{lin2023}残基レベル埋め込み（6.5億パラメータモデル、フリーズ）を糖鎖付加部位の位置情報を組み込んだsite-aware MLPを通じて射影することでエンコードされる。疾患および経路ノードは、PubMedBERTテキスト埋め込み（1.1億パラメータ、フリーズ）を学習済みMLPを通じて射影することでエンコードされる。

\textbf{リンク予測の訓練。}
モデルは型制約付き負例サンプリング（正例あたり64負例）と自己敵対的重み付けを用いてKGリンク予測タスクで訓練される。スコアリング関数はDistMultとRotatEのサブスコアラーを関係ごとの適応的重みで組み合わせる。訓練損失は二値交差エントロピーに構造的対照正則化（InfoNCE, $\tau = 0.07$）とL2正則化を加えたものである。最適化にはコサインアニーリングと混合精度訓練を伴うAdamを使用する。

\textbf{帰納的評価。}
帰納的設定で訓練を行い、エンティティの15\%（そのエッジを含む）を訓練中にホールドアウトし、未見のエンティティを含むトリプルの埋め込みとスコアリング能力を評価する。ホールドアウトセットでの帰納的フィルタードMRRは0.207（Hits@10 = 0.522）であり、モダリティ特異的エンコーダが訓練中に見ていないエンティティに汎化することを実証している。これは糖鎖予測アプリケーションにおいて不可欠であり、モデルは訓練KGに存在しない可能性のある糖鎖を埋め込む必要がある。

\subsection{階層的予測タスク}

\textbf{N結合型部位分類。}
2層MLP（256 $\to$ 128 $\to$ 1）がアスパラギン位置でのESM-2残基レベル埋め込みを入力として受け取り、N結合型糖鎖修飾確率を予測する。訓練にはクラスバランスサンプリングを伴う二値交差エントロピーを使用する。ホールドアウトタンパク質での評価ではF1、精度、再現率、AUC-ROCを報告する。

\textbf{部位ランキング。}
複数の確認された糖鎖付加部位を持つタンパク質について、ペアワイズランキングモデルがESM-2埋め込みとマージンベースのランキング損失を使用して部位を糖鎖付加傾向によりスコアリングする。評価ではタンパク質ごとの既知部位間のRecall@$k$とMRRを報告する。

\textbf{糖鎖修飾タイプ予測。}
分数回帰モデルが各部位でのハイマンノース型、複合型、ハイブリッド型、非修飾状態の比率組成をESM-2埋め込みを使用して予測する。優勢タイプはargmax予測として取得され、4クラスにわたって精度を計算する。

\textbf{構造ファミリー分類。}
8クラスMLP分類器が各糖鎖観測を糖鎖分類学から導出された構造ファミリーに割り当てる。入力特徴量はESM-2残基レベル埋め込みであり、評価ではTop-1およびTop-2精度を報告する。

\textbf{構造クラスタ予測。}
訓練済みKGモデルからの糖鎖埋め込みをKMeans（$K = 20$）でクラスタリングし、検索モデルのTop-$k$クラスタ割り当て（$k$最近傍糖鎖埋め込み間の多数決による）を真の糖鎖のクラスタに対して評価する。検索モデルは部位レベルタンパク質表現（ESM-2 MLP）と糖鎖埋め込みの間の共有256次元空間でのコサイン類似度を使用する。

\textbf{正確な糖鎖同定。}
同じ検索モデルが、クエリのタンパク質/部位表現とのコサイン類似度により全2,357のN結合型糖鎖候補をランキングする。2つのバリアントが評価される：部位レベル（糖鎖付加位置での残基レベルESM-2特徴量を使用）およびタンパク質レベル（3,345候補にわたる平均プーリングESM-2特徴量を使用）。評価はフィルタードランキングプロトコル~\cite{bordes2013}に従い、MRR、Hits@1、Hits@3、Hits@10、平均ランクを報告する。

\subsection{細胞型条件付き糖鎖予測}

\textbf{データ準備。}
細胞型酵素発現プロファイルはTabula Sapiensアトラス~\cite{tabulasapiens2022}から取得される：748細胞型と339糖鎖修飾関連酵素遺伝子のそれぞれについて、平均対数正規化発現値が記録される。glycoMusubiからの酵素-糖鎖生合成エッジは、各酵素遺伝子を生成することが知られている糖鎖構造のセットにマッピングする。

\textbf{クラスタ内酵素スコアリング。}
各細胞型$c$とクラスタ$k$内の糖鎖$g$について、酵素スコアは：
\begin{equation}
s(g, c) = \sum_{e \in \text{enzymes}(g)} \text{expr}(e, c)
\end{equation}
ここで$\text{enzymes}(g)$はKGにおいて糖鎖$g$にリンクされた酵素遺伝子のセット、$\text{expr}(e, c)$は細胞型$c$における酵素$e$の発現レベルである。酵素リンクのない糖鎖はスコアゼロとなる。

\textbf{評価モード。}
各クラスタ内で4つのモードが評価される：(1)~\textbf{周辺}：全748細胞型にわたる酵素スコアの平均；(2)~\textbf{オラクル}：真の糖鎖に最良のランクを与える細胞型（細胞型選択の上限）；(3)~\textbf{人気度}：訓練セットの糖鎖出現頻度によるランキング；(4)~\textbf{酵素+人気度}：バリデーションセットで$\alpha$を選択した線形ブレンド$s_{\text{blend}} = \alpha \cdot s_{\text{enzyme}} + (1 - \alpha) \cdot s_{\text{pop}}$。

\textbf{学習スコアラー。}
CellTypeGlycanScorerニューラルネットワークは、酵素発現ベクトル（339次元）をLayerNormとGELU活性化を持つ2層MLPを通じて、糖鎖埋め込み（KGモデルからの256次元）を並列射影を通じてマッピングし、内積（学習可能な温度と糖鎖ごとのバイアス付き）で各糖鎖をスコアリングする正規化ベクトルを生成する。訓練は各クラスタ内でInfoNCE損失を使用し、細胞型発現は訓練セット全体で平均化される（サンプルごとの真の細胞型が不明なため）。

\textbf{カスケード推定。}
フルランキングカスケードMRRは、クラスタレベルのHits@1（0.166）とクラスタ内MRRの積として推定され、エンドツーエンドの性能の下限を提供する。

\subsection{統計的手法}

すべての実験はタンパク質で層化された単一のtrain/validation/testスプリットを使用する（タンパク質が複数のスプリットに現れることはない）。検索メトリクスはフィルタードランキングプロトコルの下で計算され、既知の真の関連はランキング候補から除外される。2層境界はモデルアーキテクチャと特徴表現にわたる一貫した性能ギャップにより検証される：クラスタレベルの正確なID予測に対する優位性は、部位レベル（0.786対0.066）およびタンパク質レベル（0.786対0.118）の両モデルで持続し、モデル固有のアーティファクトを除外する。

\subsection{実装}

glycoMusubiはPyTorchとPyTorch Geometricを使用してPython 3.11で実装されている。事前計算されたESM-2残基レベル埋め込みとPubMedBERTテキスト埋め込みは\texttt{.pt}ファイルとしてキャッシュされる。設定はOmegaConfによる階層的YAMLファイルで管理される。完全なパイプライン---ETL、埋め込み、予測、評価---はスクリプト化されたエントリポイントを通じて実行可能である。再現可能なデプロイメントのためにすべての依存関係を含むDockerイメージが提供される。

% ══════════════════════════════════════════════════════════════════
\section{結果}

\subsection{ナレッジグラフの構築}

我々は自動化された4段階ETLパイプライン（ダウンロード、クリーニング、ビルド、バリデーション）を通じてglycoMusubiを構築した。このパイプラインは6つの公開データベースからのデータを統合する（図~\ref{fig:kg_benchmark}a）。得られたナレッジグラフは10の生物学的エンティティ型（糖鎖、タンパク質、酵素、疾患、バリアント、化合物、PTM部位、モチーフ、反応、細胞内局在）にわたる78,263ノードと14の関係型にわたる2,508,960エッジを含む（表~\ref{tab:stats}）。

識別子の正規化はデータベース間の競合を解決する：UniProtアイソフォームアクセッションはカノニカル形式にマッピングされ、GlyTouCanアクセッションはソースパターンに対して検証され、糖鎖付加部位は明確な相互参照を可能にする構造化識別子としてエンコードされる。

\begin{table}[ht]
\centering
\caption{glycoMusubiの統計。}
\label{tab:stats}
\small
\begin{tabular}{lrl}
\toprule
\textbf{エンティティ型} & \textbf{件数} & \textbf{主要ソース} \\
\midrule
糖鎖 & 58,628 & GlyTouCan \\
タンパク質 & 4,372 & UniProt \\
酵素 & 1,459 & GlyGen \\
部位 & 10,218 & UniProt, PTMCode \\
疾患 & 1,936 & UniProt \\
化合物 & 746 & ChEMBL \\
バリアント & 522 & UniProt \\
モチーフ & 181 & GlyTouCan \\
反応 & 164 & Reactome \\
細胞内局在 & 37 & UniProt \\
\midrule
\textbf{総ノード数} & \textbf{78,263} & \\
\textbf{総エッジ数} & \textbf{2,508,960} & 14関係型 \\
\bottomrule
\end{tabular}
\end{table}

\subsection{階層的糖鎖予測ベンチマーク}

糖鎖予測を6つのタスクの階層に分解した。各タスクは候補空間を狭め、特異性を増加させる（図~\ref{fig:kg_benchmark}b、表~\ref{tab:hierarchical}）。

\textbf{タスク1：N結合型部位分類。}
ESM-2~\cite{lin2023}残基レベル埋め込みを2層MLPへの入力として使用し、ホールドアウトタンパク質でF1 = 0.924（AUC-ROC = 0.825）を達成した。これは局所配列コンテキストがN結合型糖鎖修飾ポテンシャルの強い予測因子であることを確認する。

\textbf{タスク2：部位ランキング。}
ESM-2ベースのランキングモデルがRecall@3 = 0.683、MRR = 0.688を達成し、タンパク質配列特徴量が部位レベルの糖鎖修飾傾向を捕捉することを示す。

\textbf{タスク3：糖鎖修飾タイプ予測。}
分数回帰モデルが4カテゴリ（ハイマンノース型、複合型、ハイブリッド型、非修飾）にわたり精度 = 0.844を達成し、糖鎖修飾の構造カテゴリがタンパク質配列に部分的にエンコードされていることを実証する。

\textbf{タスク4：構造ファミリー（8クラス）。}
ESM-2 MLPがTop-2精度 = 0.858、Top-1精度 = 0.638を達成し、タンパク質配列が糖鎖を少数の構造ファミリーに確実に絞り込むことを示唆する。

\textbf{タスク5：構造クラスタ（$K = 20$）。}
クラスタはKG由来の糖鎖埋め込みに対するKMeansにより定義される（平均クラスタサイズ$\sim$118糖鎖）。部位レベル検索モデルがHits@5 = 0.786、Hits@10 = 0.978を達成し、タンパク質配列と部位コンテキストが正しい構造的近傍を高精度で予測することを示す。

\textbf{タスク6：正確な糖鎖同定。}
部位レベル検索は2,357のN結合型糖鎖に対してランキングし、MRR = 0.066（Hits@10 = 0.145）を達成した。タンパク質レベル検索は3,345糖鎖に対してランキングし、MRR = 0.118（Hits@10 = 0.228）を達成した。クラスタレベルの性能（タスク5）と正確な同定（タスク6）の間の約10倍のギャップが、根本的な決定境界を明らかにする。

\begin{table}[ht]
\centering
\caption{階層的糖鎖予測の結果。}
\label{tab:hierarchical}
\small
\begin{tabular}{llllrlr}
\toprule
\textbf{タスク} & \textbf{候補} & \textbf{モデル} & \textbf{メトリクス} & \textbf{値} & \textbf{副次メトリクス} & \textbf{値} \\
\midrule
N結合型分類 & 二値 & ESM-2 MLP & F1 & 0.924 & AUC-ROC & 0.825 \\
部位ランキング & タンパク質内 & ESM-2 ranking & Recall@3 & 0.683 & MRR & 0.688 \\
糖鎖修飾タイプ & 4クラス & Frac. regression & Accuracy & 0.844 & & \\
構造ファミリー & 8クラス & ESM-2 MLP & Top-2 Acc & 0.858 & Top-1 Acc & 0.638 \\
構造クラスタ & 20クラスタ & Site-level retr. & H@5 & 0.786 & H@10 & 0.978 \\
正確な糖鎖（部位） & 2,357糖鎖 & Site-level retr. & MRR & 0.066 & H@10 & 0.145 \\
正確な糖鎖（タンパク質） & 3,345糖鎖 & Protein-level retr. & MRR & 0.118 & H@10 & 0.228 \\
\bottomrule
\end{tabular}
\end{table}

O結合型予測問題は示唆的な対照を提供する：はるかに小さな候補空間では、部位レベル検索がMRR = 0.847を達成し（補足表~S2）、N結合型予測の困難さが本質的な予測不可能性よりも大きく多様な候補空間から主に生じることを示唆する。

\subsection{細胞型酵素発現解析}

細胞コンテキストが決定ギャップを橋渡しできるかを調査するため、配列ベースのクラスタ予測と細胞型特異的酵素発現スコアリングを組み合わせたカスケードモデルを開発した（図~\ref{fig:twolayer}a）。カスケードは糖鎖予測を以下のように分解する：
\begin{equation}
P(\text{glycan} \mid \text{protein, site, cell\_type}) = P(\text{cluster} \mid \text{protein, site}) \times P(\text{glycan} \mid \text{cluster, cell\_type})
\end{equation}

Tabula Sapiensシングルセルアトラス~\cite{tabulasapiens2022}からの酵素発現プロファイル（748細胞型、339糖鎖修飾関連酵素遺伝子）とglycoMusubiからの酵素-糖鎖生合成エッジを使用し、予測された各クラスタ内で生産酵素の発現レベルにより糖鎖をスコアリングした。

結果は重大なデータギャップを明らかにする（表~\ref{tab:celltype}）。クラスタコンテキスト内（平均$\sim$118候補）では、人気度ベースラインがMRR = 0.105を達成する一方、酵素発現周辺はMRR = 0.035、オラクルはMRR = 0.026にとどまる。酵素+人気度ブレンドは酵素成分に最適重み$\alpha = 0.0$を割り当てる。学習ニューラルスコアラーはMRR = 0.104を達成し---人気度ベースラインと同等である。

根本原因はアノテーションの疎さである：2,357のN結合型糖鎖のうちわずか929（39.4\%）がglycoMusubiにおいて何らかの酵素リンクを持ち、これらは1,843糖鎖（全N結合型糖鎖の78\%）を含む単一のメガクラスタに集中している。20の構造クラスタのうち、非自明な酵素カバレッジを持つのは1つのみである。

\begin{table}[ht]
\centering
\caption{細胞型条件付き糖鎖予測。}
\label{tab:celltype}
\small
\begin{tabular}{llrrrp{3.5cm}}
\toprule
\textbf{手法} & \textbf{スコープ} & \textbf{MRR} & \textbf{H@10} & \textbf{H@50} & \textbf{備考} \\
\midrule
人気度ベースライン & 全体(2,357) & 0.100 & 0.193 & 0.439 & 訓練頻度 \\
酵素周辺 & 全体(2,357) & 0.001 & 0.000 & 0.000 & 全細胞型平均 \\
\midrule
人気度 & クラスタ内 & 0.105 & 0.207 & 0.469 & 平均$\sim$118候補 \\
酵素周辺 & クラスタ内 & 0.035 & 0.034 & 0.034 & 平均$\sim$118候補 \\
酵素オラクル & クラスタ内 & 0.026 & 0.026 & 0.029 & サンプルごと最良細胞型 \\
酵素+人気度 & クラスタ内 & 0.105 & 0.207 & 0.469 & 最適$\alpha = 0.0$ \\
学習スコアラー & クラスタ内 & 0.104 & 0.205 & 0.467 & InfoNCE訓練 \\
\midrule
部位レベル検索 & 全体(2,357) & 0.066 & 0.145 & --- & 配列のみ（参照） \\
タンパク質レベル検索 & 全体(3,345) & 0.118 & 0.228 & --- & 配列のみ（参照） \\
\bottomrule
\end{tabular}
\end{table}

\subsection{体系的な陰性結果}

いくつかの追加アプローチもベースラインを超える予測改善に失敗した（表~\ref{tab:negative}）。タンパク質配列からの単糖組成予測はすべての糖種で$R^2 < 0$であり、学習可能な関係がないことを確認した。生合成経路データベース（glycoworkサブグラフ、WURCSベースの反応）は人気度ベースラインを改善しなかった。Human Protein Atlas~\cite{uhlen2015}からのRNA発現レベルはMRR = 0.001を達成し、人気度ベースラインを大幅に下回った。検索補助シグナルとしての酵素経路スコアリングは、検索スコア単独に対する最適ブレンド重み$\alpha = 1.0$をもたらした。

\begin{table}[ht]
\centering
\caption{陰性結果のまとめ。}
\label{tab:negative}
\small
\begin{tabular}{llcrp{3.5cm}}
\toprule
\textbf{アプローチ} & \textbf{データソース} & \textbf{メトリクス} & \textbf{値} & \textbf{解釈} \\
\midrule
組成予測 & ESM-2 $\to$ 単糖 & $R^2$ & $< 0$ & 配列から学習不可 \\
反応ネットワーク & glycowork & MRR & $\leq$ pop. & アノテーション不足 \\
組織発現 & HPA RNA consensus & MRR & 0.001 & バルク組織では不十分 \\
酵素経路 & GlyGen酵素エッジ & $\alpha^*$ & 1.0 & 追加情報なし \\
細胞型酵素 & Tabula Sapiens & MRR & 0.035 & アノテーション不足 \\
\bottomrule
\end{tabular}
\end{table}

\subsection{2層決定モデル}

階層的ベンチマークと陰性結果の解析は合わせて、2つの異なる層を持つ糖鎖修飾決定の定量的モデルを支持する（図~\ref{fig:twolayer}b）。

\textbf{第1層：配列決定。}
タンパク質配列と部位コンテキストは糖鎖修飾の構造クラスを強く制約する（タスク1--5; 部位分類でF1 = 0.924、構造クラスタでH@5 = 0.786）。生物学的基盤は糖転移酵素の基質特異性である~\cite{schjoldager2020,narimatsu2019}。

\textbf{第2層：環境決定。}
構造クラス内では、正確な糖鎖のアイデンティティは細胞因子に依存する：酵素発現、ヌクレオチド糖の利用可能性、ゴルジ体通過時間、競合する生合成経路~\cite{stanley2022,varki2022essentials}。これはタスク6に対応し、性能が急激に低下する（MRR = 0.066--0.118）。我々の細胞型解析は、現在のアノテーションがこのギャップを計算的に橋渡しするには疎すぎることを示す（39\%の糖鎖カバレッジ、1クラスタに集中）。

この境界は使用したモデルの限界ではない：10倍のギャップはモデルアーキテクチャ（部位レベル対タンパク質レベル検索）および特徴表現（ESM-2 MLP対埋め込みベースの検索）にわたって持続し、糖鎖修飾生物学の根本的な特性を反映していることを示唆する（図~\ref{fig:twolayer}c）。

% ══════════════════════════════════════════════════════════════════
\section{考察}

我々の階層的ベンチマークは、タンパク質配列が糖鎖修飾について何を予測でき何を予測できないかの最初の体系的定量化を提供する。構造クラス予測の高精度（$K = 20$クラスタでH@5 = 0.786）は、糖転移酵素の特定のアクセプター配列に対する既知の特異性~\cite{schjoldager2020}と整合し、配列ベースのツールが所与の部位で\textit{どのタイプ}の糖鎖修飾が起こるかを確実に予測できることを示唆する。このレベルの予測は構造生物学、分子動力学シミュレーション、合理的タンパク質工学の多くの応用にとって十分である。

正確な糖鎖同定でMRR = 0.066--0.118に急落することは、正確な糖鎖構造の予測にはタンパク質配列以外の情報が必要であることを意味する。候補情報源---シングルセル酵素発現、組織発現プロファイル、反応ネットワーク、組成---の体系的調査では、現在このギャップを橋渡しするものは見つからなかった。

酵素発現の失敗は特に示唆的である。Tabula Sapiensアトラスは利用可能な最も詳細な細胞型特異的酵素発現の見方を提供するが、予測シグナルを追加しない。これは生物学的仮説が間違っているからではない---糖鎖生合成は実際に酵素発現によって決定される---しかし、酵素-糖鎖アノテーションマッピングが決定的に不完全だからである。現在のデータベースにおいて2,357のN結合型糖鎖のうちわずか929（39.4\%）が何らかの酵素リンクを持ち、単一の構造クラスタに集中している。このアノテーションギャップが、モデル能力ではなく、主要なボトルネックである。

O結合型予測結果（MRR = 0.847）は重要な文脈を提供する。O結合型糖鎖---候補空間が小さい---に対する高精度は、N結合型の課題が本質的な予測不可能性よりもむしろ候補空間のサイズと糖鎖の多様性の問題であることを示唆する。

\textbf{先行研究との比較。}
我々の研究はInSaNNE~\cite{kellman2024}と最も直接的に比較可能であり、InSaNNEはリカレントニューラルネットワークを使用してタンパク質配列特徴量から糖鎖構造を予測する。InSaNNEは199糖鎖と57糖鎖付加部位にわたる二値有無分類でAUC = 0.92を達成した。我々のアプローチは3つの点で異なる：(i) 199に対する二値分類ではなく2,357糖鎖にわたる検索ランキング（MRR）を評価し、問題の完全なスケールを反映する；(ii) 問題を6つのタスクの階層に分解し、予測がどこで成功しどこで失敗するかを明らかにする；(iii) 追加情報源としての細胞コンテキストを体系的に調査する。我々が同定した決定境界---確実な構造クラス予測だが不良な正確ID予測---はGlycoimpactの「制約付き生合成」フレームワーク~\cite{glycoimpact2024}と整合し、タンパク質配列が糖鎖を制約するが決定はしないことを示した。

Yadavら~\cite{yadav2022}はゴルジ体糖鎖プロセシングを最適情報符号化チャネルとしてモデル化し、構造多様性が酵素反応速度論と区画組織化によって制約されると予測した。我々の経験的決定境界はこの理論的枠組みに定量的支持を提供する：クラスタレベルと正確なID予測の間の10倍のギャップは、ゴルジ体における区別可能な酵素状態の数によって課される情報理論的限界と整合する。

Chrysinasら~\cite{chrysinas2024}は同じTabula Sapiensアトラスを使用して糖鎖修飾酵素の発現を解析し、末端プロセシング酵素（例：シアリルトランスフェラーゼ、フコシルトランスフェラーゼ）が高度に細胞型特異的である一方、コア経路酵素（例：MGAT1、MGAT2）はユビキタスに発現していることを見出した。彼らの定性的観察は我々の定量的知見と一致する：細胞型で変動する酵素はまさに最も疎な糖鎖アノテーションを持つもの（39\%カバレッジ）であり、それらの判別的使用を妨げている。

SweetNet~\cite{burkholz2022}、GIFFLAR~\cite{thomes2024}、GlycanML~\cite{glyml2024}は構造から糖鎖特性を予測する（順方向）が、逆問題には取り組んでいない。CoNglyPred~\cite{congly2025}はESM-2とAlphaFold2の特徴量を使用してN結合型糖鎖付加部位を予測し、我々のタスク1（F1 = 0.924）を補完するが、糖鎖のアイデンティティには取り組んでいない。glycoMusubiリソースは先行アプローチが個別に使用していた6つのデータベースからのデータを統合し、単一データベースアプローチでは不可能なマルチホップクエリを可能にする。

\textbf{限界。}
ベンチマークはGlyConnectの部位特異的糖鎖割り当てに基づいており、よく研究されたタンパク質に偏っている。構造クラスタリングはKG由来の糖鎖埋め込みに依存する。細胞型解析はRNAレベルでの酵素発現が酵素活性を反映すると仮定している。現在のKGには質量分析による糖鎖同定データは含まれていない。

\textbf{今後の方向性。}
決定ギャップを埋めるにはデータと手法の並行的な取り組みが必要である。酵素-糖鎖の生合成関係の体系的解明---特に複合型N結合型糖鎖について---はカバレッジを直接的に増加させる。空間グライコミクスと質量分析ベースの糖鎖同定の統合は、正確な糖鎖レベルでの予測を大幅に改善し得る。階層的ベンチマークは進歩を測定するための標準化された評価フレームワークを提供する。

% ══════════════════════════════════════════════════════════════════
\section*{データの利用可能性}

glycoMusubiは以下の公開データを統合する：GlyGen（\url{https://data.glygen.org}）、GlyTouCan（\url{https://glytoucan.org}）、UniProt（\url{https://www.uniprot.org}）、ChEMBL（\url{https://www.ebi.ac.uk/chembl}）、PTMCode（\url{https://ptmcode.embl.de}）、GlyConnect（\url{https://glyconnect.expasy.org}）。構築されたナレッジグラフ、学習済み埋め込み、ベンチマーク結果は出版時にZenodoに登録予定である。

\section*{資金}

【追記予定】

% ══════════════════════════════════════════════════════════════════
% Figures
% ══════════════════════════════════════════════════════════════════

\begin{figure}[!ht]
\centering
\includegraphics[width=\textwidth]{../../experiments_v2/paper_figures/fig1_kg_benchmark.pdf}
\caption{\textbf{glycoMusubiリソースと階層的ベンチマークの概要。}
(a)~データ統合パイプライン：6つの公開データベースが自動化された4段階ETLパイプラインを通じて処理され、78,263ノード（10型）と2,508,960エッジ（14関係型）を持つヘテロジニアスナレッジグラフが構築される。(b)~階層的予測ベンチマーク：特異性が増加する5つのタスクにわたる性能。破線は構造クラス予測から正確な糖鎖同定へと性能が急落する決定境界を示す。}
\label{fig:kg_benchmark}
\end{figure}

\begin{figure}[!ht]
\centering
\includegraphics[width=\textwidth]{../../experiments_v2/paper_figures/fig2_twolayer.pdf}
\caption{\textbf{2層決定モデルと細胞型解析。}
(a)~配列ベースのクラスタ予測（第1層）と細胞型酵素発現スコアリング（第2層）を組み合わせたカスケードアーキテクチャ。(b)~糖鎖修飾決定の概念モデル：タンパク質配列が構造クラスを決定し（第1層; H@5 = 0.786）、正確な糖鎖のアイデンティティには細胞コンテキストが必要（第2層; MRR = 0.066--0.118）。10倍のギャップはモデルアーキテクチャにわたって持続する。(c)~予測階層にわたる性能。青いバーは配列決定タスク（第1層）、赤いバーは環境依存タスク（第2層）を表す。}
\label{fig:twolayer}
\end{figure}

% ══════════════════════════════════════════════════════════════════
% References
% ══════════════════════════════════════════════════════════════════

\begin{thebibliography}{27}

\bibitem{varki2017}
Varki,A. (2017) Biological roles of glycans. \textit{Glycobiology}, \textbf{27}, 3--49.

\bibitem{reily2019}
Reily,C. \textit{et al.} (2019) Glycosylation in health and disease. \textit{Nat. Rev. Nephrol.}, \textbf{15}, 346--366.

\bibitem{schjoldager2020}
Schjoldager,K.T. \textit{et al.} (2020) Global view of human protein glycosylation pathways and functions. \textit{Nat. Rev. Mol. Cell Biol.}, \textbf{21}, 729--749.

\bibitem{narimatsu2019}
Narimatsu,Y. \textit{et al.} (2019) An atlas of human glycosylation pathways enables display of the human glycome by gene engineering. \textit{Mol. Cell}, \textbf{75}, 394--407.

\bibitem{stanley2022}
Stanley,P. \textit{et al.} (2022) N-Glycans. In: Varki,A. \textit{et al.} (eds) \textit{Essentials of Glycobiology}, 4th edn, Ch.~9. Cold Spring Harbor Laboratory Press.

\bibitem{varki2022essentials}
Varki,A. \textit{et al.} (2022) \textit{Essentials of Glycobiology}, 4th edn. Cold Spring Harbor Laboratory Press.

\bibitem{tiemeyer2017}
Tiemeyer,M. \textit{et al.} (2017) GlyTouCan: an accessible glycan structure repository. \textit{Glycobiology}, \textbf{27}, 915--919.

\bibitem{uniprot2023}
The UniProt Consortium (2023) UniProt: the Universal Protein Knowledgebase in 2023. \textit{Nucleic Acids Res.}, \textbf{51}, D523--D531.

\bibitem{york2020}
York,W.S. \textit{et al.} (2020) GlyGen: computational and informatics resources for glycoscience. \textit{Glycobiology}, \textbf{30}, 72--73.

\bibitem{zdrazil2024}
Zdrazil,B. \textit{et al.} (2024) The ChEMBL Database in 2023. \textit{Nucleic Acids Res.}, \textbf{52}, D1180--D1192.

\bibitem{alocci2019}
Alocci,D. \textit{et al.} (2019) GlyConnect: glycoproteomics goes visual, interactive, and analytical. \textit{J. Proteome Res.}, \textbf{18}, 698--708.

\bibitem{burkholz2022}
Burkholz,R. \textit{et al.} (2022) SweetNet: a neural network for predicting glycan properties from graph representation. \textit{Bioinformatics}, \textbf{38}, 4599--4605.

\bibitem{thomes2024}
Thom\`es,L. and Burkholz,R. (2024) GIFFLAR: a higher-order graph neural network for glycan function prediction. \textit{Preprint at bioRxiv}.

\bibitem{thomes2021}
Thom\`es,L. \textit{et al.} (2021) glycowork: a Python package for glycan data science and machine learning. \textit{Glycobiology}, \textbf{31}, 1240--1244.

\bibitem{lin2023}
Lin,Z. \textit{et al.} (2023) Evolutionary-scale prediction of atomic-level protein structure with a language model. \textit{Science}, \textbf{379}, 1123--1130.

\bibitem{bordes2013}
Bordes,A. \textit{et al.} (2013) Translating embeddings for modeling multi-relational data. \textit{Adv. Neural Inf. Process. Syst.}, \textbf{26}.

\bibitem{tabulasapiens2022}
The Tabula Sapiens Consortium (2022) The Tabula Sapiens: a multiple-organ, single-cell transcriptomic atlas of humans. \textit{Science}, \textbf{376}, eabl4896.

\bibitem{uhlen2015}
Uhl\'en,M. \textit{et al.} (2015) Tissue-based map of the human proteome. \textit{Science}, \textbf{347}, 1260419.

\bibitem{minguez2015}
Minguez,P. \textit{et al.} (2015) PTMcode v2: a resource for functional associations of post-translational modifications within and between proteins. \textit{Nucleic Acids Res.}, \textbf{43}, D503--D509.

\bibitem{kellman2024}
Kellman,B.P. \textit{et al.} (2024) Elucidating Human Protein Glycosylation using an Informed Search of Artificial Neural Network Ensembles (InSaNNE). \textit{Preprint at bioRxiv}, doi:10.1101/2024.02.02.578654.

\bibitem{bojar2022}
Bojar,D. and Lisacek,F. (2022) Glycoinformatics in the Artificial Intelligence Era. \textit{Chem.\ Rev.}, \textbf{122}, 15971--15988.

\bibitem{congly2025}
Li,M. \textit{et al.} (2025) CoNglyPred: accurate prediction of protein N-linked glycosylation sites combining ESM-2 and AlphaFold2 features. \textit{PROTEOMICS}, \textbf{25}, e2400129.

\bibitem{glyml2024}
Bojar,D. \textit{et al.} (2024) GlycanML: A Multi-Task and Multi-Structure Benchmark for Glycan Machine Learning. \textit{Preprint at arXiv}, arXiv:2405.16206.

\bibitem{yadav2022}
Yadav,V. \textit{et al.} (2022) Golgi glycan processing as an optimal information coding channel. \textit{eLife}, \textbf{11}, e76505.

\bibitem{chrysinas2024}
Chrysinas,P. \textit{et al.} (2024) Single cell analysis of human glycosyltransferases and glycosidases. \textit{NAR Genomics Bioinformatics}, \textbf{6}, lqae069.

\bibitem{glycoimpact2024}
Lundstr\"om,J. \textit{et al.} (2024) GlycoImpact: In-silico glycan scoring by amino acid substitution matrix. \textit{Preprint at bioRxiv}, doi:10.1101/2024.09.16.613237.

\end{thebibliography}

\end{document}
